\chapter*{Abstract}
This project presents a local-first Retrieval-Augmented Generation (RAG) system designed for querying publicly available web content from Southern University of Science and Technology (SUSTech). The system integrates web data acquisition, text preprocessing, semantic chunking, vector-based retrieval, reranking, and large language model generation into a unified pipeline for natural language question answering.

The backend is implemented in Python with FastAPI, supporting configurable LLM backends including llama.cpp and vLLM. It adopts BGE models for embedding and reranking, and Qwen3-based models for response generation. The frontend is developed using Vue 3 and Vite, providing a responsive web interface with streaming responses via Server-Sent Events (SSE), multi-session support, and customizable UI themes.

The system is designed to be modular and deployment-flexible, enabling operation across local environments and GPU-enabled servers. By combining efficient retrieval mechanisms with generative models, the system significantly improves the accessibility and usability of campus information.